\documentclass[12pt,a4paper]{article}

%% Formatting per IRCFGS 2026 author guidelines:
%% Times New Roman 12pt, 1.15 spacing, justified, margins T/B 1in, L/R 0.8in.
\usepackage[T1]{fontenc}
\usepackage{newtxtext,newtxmath}   % Times-like text and math
\usepackage[a4paper,top=1in,bottom=1in,left=0.8in,right=0.8in]{geometry}
\usepackage{setspace}
\setstretch{1.15}
\usepackage{amssymb}
\usepackage[hidelinks]{hyperref}
\setlength{\parindent}{0pt}
\setlength{\parskip}{6pt}

\begin{document}

\begin{center}
{\large\bfseries Adversarial robustness and privacy-utility tradeoffs in
real-time graph-based fraud ring detection}\\[8pt]
O.~Karunaratne\textsuperscript{1*} and T.~Kaushalya\textsuperscript{1}\\[4pt]
\textit{\textsuperscript{1}Affiliation to be confirmed}\\
\texttt{karunaratneovindu@gmail.com}\\[2pt]
{\small\textit{(Draft for review --- IRCFGS 2026, Faculty of Graduate Studies,
University of Colombo)}}
\end{center}

\vspace{6pt}
\textbf{Keywords:} fraud detection; weakly connected components; adversarial
robustness; differential privacy; streaming systems

%% ======================================================================
\section*{Abstract}

Weakly connected component (WCC) graph analysis is the dominant technique for
real-time fraud ring detection in streaming financial systems, yet its
robustness to adversarial evasion in production, multi-relational deployments
remains largely unstudied. We present a streaming fraud-detection testbed
(Apache Flink, Memgraph, LangGraph) and use it to answer two questions: how WCC
ring detection degrades when an attacker rotates connections across three
independent channels (devices, IP addresses, merchants), and what
privacy-utility tradeoff arises when Laplace differential privacy is applied to
the co-located velocity detector. Across 30 Monte Carlo trials per
configuration, we find that WCC exhibits strong multi-channel redundancy:
severing only device and IP connections, or only merchant connections, leaves
recall unchanged ($R = 1.000$). Only simultaneous evasion of all three channels
degrades detection, and an attacker must fully evade roughly 60--80\% of ring
members before recall collapses. We show this transition is not incidental but
follows a closed-form law: a ring survives precisely when at least half its
members remain in the shared component (the Intersection-over-Union matching
threshold), and the aggregate recall curve is predicted, without fitting, by a
hypergeometric coverage model that matches our empirical results to three
decimal places. For differential privacy, we identify a privacy budget of
$\varepsilon \approx 5$ as the practical operating point: below it,
false-positive rates become operationally unacceptable (FPR $=0.28$ at
$\varepsilon=1$); above it, near-perfect recall is preserved (FPR $=0.027$ at
$\varepsilon=5$). Our results quantify the adversarial cost of evading
production WCC pipelines and expose merchant-side edges as the highest-cost
channel to corrupt. The policy implication is concrete: graph-integrity
controls should concentrate on merchant-side relationships, and
privacy-preserving fraud statistics are operationally viable only above a
quantifiable privacy budget.

%% ======================================================================
\section*{Extended abstract}

\subsection*{A brief introduction}

Fraud rings --- coordinated groups of accounts that share devices, IP
addresses, and merchant relationships --- account for a disproportionate share
of financial-fraud losses. Graph-based detection using weakly connected
components (WCC) is among the most widely deployed production techniques, owing
to its simplicity and interpretability, and underpins systems such as AWS Graph
Fraud Detection and Neo4j Fraud Detection. Despite this ubiquity, two questions
remain underexplored. \textbf{First (adversarial robustness):} how many
adversarial modifications must an attacker make to evade WCC-based detection,
and which connection channels are the weakest links? Prior adversarial work on
graphs targets learned classifiers or community detection on static,
single-relation social graphs; the streaming, multi-relational, rule-based WCC
deployments characteristic of production fraud pipelines --- where an attacker
acts transaction-by-transaction and must simultaneously evade redundant
channels --- remain unstudied. \textbf{Second (privacy-utility tradeoff):} if
an institution publishes aggregate velocity statistics under differential
privacy, what privacy budget is required before noise degrades operational
detection? Our objective is to answer both empirically and to explain the
adversarial result with a predictive theoretical model.

\subsection*{Main body}

\textbf{Methodology.} We build a streaming fraud-detection testbed (Apache
Flink, Memgraph, LangGraph) and a reproducible simulation harness that
replicates a production WCC ring detector over a multi-relational property graph
with four node types (User, Merchant, Device, IP). A ring of $k$ users sharing
devices, IPs, and merchants forms a connected component; components larger than
three nodes are flagged. We adopt a white-box threat model in which the attacker
knows the algorithm, thresholds, and ring structure, and is given an
adversarial budget $B$ equal to the number of ring users whose transaction
parameters may be modified. We evaluate five channel-aware strategies:
device/IP rotation, merchant rotation, full evasion (all three channels), a
complementary velocity-throttle attack, and their combination, across the full
budget range. Detection is scored against ground-truth rings using
Intersection-over-Union (IoU) matching at 0.5, over 30 independent Monte Carlo
trials per configuration. We additionally apply the Laplace mechanism to the
velocity detector's count and sum statistics, sweeping the privacy budget
$\varepsilon$ across $\{0.1 \ldots 50\}$, and we evaluate an optimal adversary
that concentrates evasions within rings, a temporal-window detector that runs
WCC over each user's recent transactions, an improved differential-privacy
mechanism, and a learned graph neural network (GCN) baseline.

\textbf{Key findings.} (1)~\textit{WCC exhibits multi-channel redundancy.}
Rotating only devices and IPs never reduces recall below 1.000 --- rings stay
connected through shared merchants --- and rotating only merchants is equally
ineffective ($F_1 \geq 0.989$), because device/IP sharing alone sustains
connectivity. (2)~\textit{The degradation follows a predictive law.}
Full-evasion recall falls sharply between 60\% and 80\% budget. We show this is
exactly the IoU matching threshold expressed per ring: a ring is detected iff at
least $\lceil \theta k \rceil$ of its $k$ members remain in the shared component
($\theta = 0.5$). Because the budget is distributed across rings, aggregate
recall equals the hypergeometric tail $P(X \leq k - \lceil \theta k \rceil)$.
This closed-form model, with no fitted parameters, reproduces the empirical
recall curve to three decimal places (predicted 0.310 vs.\ measured 0.313 at
60\% budget). (3)~\textit{The velocity-throttle attack is orthogonal,} defeating
the velocity detector without affecting WCC ($F_1 = 0.998$). (4)~\textit{No
single-snapshot defense recovers the ring} once all channels are severed,
indicating the vulnerability is structural. (5)~\textit{For differential
privacy,} FPR drops from 0.28 at $\varepsilon=1$ to 0.027 at $\varepsilon=5$; a
count-primary mechanism that avoids the high-sensitivity sum channel attains the
same false-positive rate at $\varepsilon=2$, a $2.5\times$ tighter budget.
(6)~\textit{Attacks are provably efficient:} the minimum cut to isolate a user
is exactly three channel rotations, and an optimal adversary that concentrates
the kill-threshold number of evasions per ring defeats every ring at 60\%
budget --- where the naive attacker still detects 31\% --- so robustness
benchmarked against naive attackers is overstated by roughly 40 budget points.
(7)~\textit{A temporal-window detector recovers fully-evaded rings:} because
full evasion is a one-time rotation, running WCC over each user's recent history
restores recall from 0 to 1.0 once the window exceeds the attacker's
sustained-evasion duration --- the first defense to beat full evasion. A learned
GCN is similarly only partially robust, retaining recall through non-structural
features rather than graph structure, reinforcing that durable defense needs
signals beyond the static graph. The closed-form model holds at $N$ up to 1000
and across ring sizes, and admits an exact deterministic robustness certificate.

\textbf{Discussion and implications.} The adversarial cost of evasion is
dominated by the hardest channel to corrupt: device rotation requires new
hardware or emulators; IP rotation is cheap via proxies; merchant rotation
requires colluding shell merchants and is the costliest. Institutions should
therefore prioritise the integrity of merchant-side edges. The coverage model
turns an empirical curve into a planning tool: given a ring-size distribution
and matching threshold, an operator can predict the evasion budget at which
detection fails. Because the failure is structural, durable defenses must move
beyond single-snapshot WCC --- which we demonstrate with a temporal-window
detector --- or incorporate non-graph signals.

\textbf{Conclusion.} We present, to our knowledge, the first
adversarial-robustness study of WCC fraud-ring detection in a streaming,
multi-relational setting, and the first to give it a predictive closed-form
model. WCC is surprisingly robust through multi-channel redundancy, but fails
structurally once all channels are severed --- at a budget our hypergeometric
model predicts without fitting. We quantify the privacy-utility tradeoff of
differentially private velocity detection and identify $\varepsilon \approx 5$
as the operating point. The policy implication for regulators and institutions
is concrete: graph-integrity controls should concentrate on merchant-side
relationships, and privacy-preserving fraud statistics are operationally viable
only above a quantifiable privacy budget.

\subsection*{Selected references}

\begin{enumerate}\setlength{\itemsep}{0pt}
  \item Akoglu, L., Tong, H., Koutra, D. (2015). Graph based anomaly detection
        and description: a survey. \textit{Data Mining and Knowledge Discovery.}
  \item Pourhabibi, T., et al. (2020). Fraud detection: A systematic literature
        review of graph-based anomaly detection approaches.
        \textit{Decision Support Systems.}
  \item Li, J., et al. (2020). Adversarial attack on community detection by
        hiding individuals. \textit{WWW.}
  \item Jia, J., et al. (2020). Certified robustness of community detection
        against adversarial structural perturbation via randomized smoothing.
        \textit{WWW.}
  \item Dwork, C., Roth, A. (2014). The algorithmic foundations of differential
        privacy. \textit{Foundations and Trends in Theoretical Computer Science.}
  \item Nissim, K., Raskhodnikova, S., Smith, A. (2007). Smooth sensitivity and
        sampling in private data analysis. \textit{STOC.}
\end{enumerate}

\end{document}
